% Generated by tools/edn_latex.py from reports/edn.md.
% Do not edit: change reports/edn.md and run `make edn`.
\ednentry{
  id = {EDN-08},
  title = {Model loading: bake into the API image via \texorpdfstring{\texttt{dvc pull}}{dvc pull} at CI build time, not the MLflow registry at runtime},
  date = {2026-09-22},
  milestone = {M4: Model Deployment},
  topic = {ML System Design},
  participants = {@lukas2510},
  decision = {The API image is built by CI with the current production model already baked in, fetched with \texttt{dvc pull} pinned to whatever \texttt{dvc.\allowbreak{}lock} on \texttt{main} points to. Promoting a model is a normal merge to \texttt{main} that updates that pointer. The API never calls the MLflow registry at build or run time; MLflow remains the experiment-tracking and audit record of which run was chosen as production.},
  rationale = {\emph{Alternatives considered:}
\begin{itemize}[leftmargin=*, topsep=2pt]
\item \textbf{Option A: pull from the MLflow model registry at API startup.} As originally stated in requirements.md FR-12.
\newline Pros: matches the target architecture's ``MLflow tracking + model registry'' component as originally drafted; a model can be promoted without a redeploy.
\newline Cons: startup now depends on DagsHub being reachable, which works against NFR-05 (ready within 30 s) and NFR-12 (a self-contained stack); a DagsHub outage right before a presentation could break the whole demo.
\item \textbf{Option B: CD resolves the MLflow registry alias and bakes the model into the image on promotion.}
\newline Pros: avoids the runtime dependency of option A while still using the registry as the deployment-time source of truth.
\newline Cons: needs a new trigger connecting an MLflow promotion event to a GitHub Actions build (webhook or poller), which is infrastructure the course does not demo yet and is disproportionate for a 5-person, one-semester, single-VM project; nothing watches it if it silently stops firing.
\item \textbf{Option C (chosen): bake into the image via \texttt{dvc pull} at CI build time; promotion is a merge to \texttt{main}.}
\newline Pros: matches the course demo's own deployment recipe (\texttt{dvc pull} the model before running \texttt{uvicorn}), just containerized; reuses GitHub Flow, already the project's branching model, as the promotion mechanism with no new trigger to build; leans on DVC, the single heaviest-weighted practice in the course rubric (15/100), for the thing that most needs to be reproducible; avoids the runtime dependency of option A.
\newline Cons: MLflow's model registry is no longer part of the deployment path, only tracking/audit, which is a smaller role than the target architecture originally sketched.
\end{itemize}
\par
\emph{Why:} Given the course's own demo pattern, the rubric's weight on DVC over MLflow, and the team's size and one-semester timeline, option C gives the same startup-independence as option B without inventing new promotion infrastructure. MLflow keeps its full credit for the experiment-tracking practice; it is just not on the deployment-time critical path.},
  ai = {Information seeking, Alternative generation, Alternative assessment, Recommendation},
  response = {Accepted},
  assessment = {While reviewing PR \#19, AI first flagged that FR-12's runtime registry pull (option A) worked against NFR-05 and NFR-12 and recommended option B. Asked to weigh that against the project's actual scope and tools, AI checked \texttt{references/\allowbreak{}course-\allowbreak{}demos.\allowbreak{}md}, found the demo's own recipe pulls the model via DVC rather than MLflow, cross-referenced the rubric's point weights, and revised its recommendation to option C, naming the trade-off (MLflow registry no longer in the deployment path) explicitly. Lukas reviewed the reasoning and confirmed option C.},
  aievidence = {Claude Code session on 2026-09-22 while reviewing PR \#19: AI recommended option B, Lukas asked ``considering our project scope and the tools we will use is this the best option for us?'', AI re-derived the recommendation from \texttt{references/\allowbreak{}course-\allowbreak{}demos.\allowbreak{}md} and the rubric weights and proposed option C instead; Lukas replied ``yes please do that''.},
  evidence = {\href{https://github.com/mlops-2627q1-mds-upc/MLOPS_Recommenditos/blob/main/docs/docs/specification.md}{Specification} FR-12; \href{https://github.com/mlops-2627q1-mds-upc/MLOPS_Recommenditos/blob/main/docs/docs/project-brief.md}{project brief} section 5 (target architecture); \href{https://github.com/mlops-2627q1-mds-upc/MLOPS_Recommenditos/blob/main/references/course-demos.md}{references/course-demos.md} ``API and deployment (M4)''; \href{https://github.com/mlops-2627q1-mds-upc/MLOPS_Recommenditos/pull/19}{PR \#19}.},
}
